\subsection{Data Structures \& Type Definitions}
\label{sec:impl_i2c_data_structures}

The I2C controller model comprises three files: \texttt{stm32f4xx\_i2c.c}, \\ \texttt{stm32f4xx\_i2c.h}, and \texttt{stm32f4xx\_types\_i2c.h}. The implementation 
file (\texttt{.c}) contains all functions that define the controller's behavior and 
enable QEMU to instantiate it. The header files (\texttt{.h}) define the C data 
structures that compose the controller model, encapsulating both the visible hardware 
state (registers accessible via MMIO) and the internal logical state required to 
emulate the I2C protocol sequence.

\subsubsection{Core Data Structure: \texttt{STM32F4XXI2CState}}

The principal structure is \texttt{STM32F4XXI2CState}, declared using QOM API macros 
to represent a complete I2C controller instance. The Code~\ref{code:i2c_state} shows 
its declaration.

\lstinputlisting[
  language=C,
  caption={\texttt{STM32F4XXI2CState} structure definition},
  label=code:{STM32F4XXI2CState}
]{Resources/Code/STM32F4XXI2CState.c}


This structure is composed of several elements that fulfill various functions. The most important ones will be analyzed individually or mentioned as the section progresses. Although not all, some of the elements that make up this structure are:

\begin{itemize}
    \item \textbf{Ten 32-bit registers} (\texttt{i2c\_cr1} through \texttt{i2c\_fltr}) 
    tahat mirror the I2C registers in the STM32F4 reference manual.
    
    \item An \texttt{I2CBus *bus} pointer enabling connection to QEMU's generic I2C 
    subsystem, allowing communication with any connected slave device model.
    
    \item Dual interrupt lines (\texttt{irq\_event} and \texttt{irq\_error}) that 
    replicate hardware separation between event-driven and error-driven interrupt 
    sources.
\end{itemize}

\subsubsection{State Tracking Structure: \texttt{stm32f4xx\_i2c\_config\_t}}

A secondary structure, \texttt{config} (type \texttt{stm32f4xx\_i2c\_config\_t}), 
tracks logic not directly visible in hardware registers (see Appendix~\ref{ann:stm32f4xx_i2c_code}). 
This structure is essential due to the STM32F405's complex flag-clearing semantics. 
For example, the \texttt{ADDR} bit in \texttt{SR1} is cleared only when \texttt{SR1} 
is read \emph{followed} by a read from \texttt{SR2}. The boolean member 
\texttt{flg\_addr} records whether \texttt{SR1} has been read, allowing the subsequent 
read from \texttt{SR2} to correctly clear the \texttt{ADDR} bit.

\subsubsection{Finite State Machine: Protocol Behavior}

The sequential nature of the I2C protocol---where register accesses trigger 
state-dependent bus operations---necessitates a finite state machine (FSM) model. 
This FSM serves as the operational heart of the emulated I2C controller, translating 
register accesses into I2C bus transactions.

The implementation defines two FSM-related types:

\begin{itemize}
    \item \textbf{\texttt{stm32f4xx\_i2c\_fsm\_trans\_t}}: A structure representing 
    a single state transition.
    
    \item \textbf{\texttt{stm32f4xx\_i2c\_fsm\_t}}: A structure containing a pointer 
    to an array of transitions and the current FSM state.
\end{itemize}

The \texttt{STM32F4XXI2CState} structure incorporates a pointer to 
\texttt{stm32f4xx\_i2c\_fsm\_t}. The FSM states, encoded as an enumeration, reflect 
the operational phases of the I2C controller in master mode as documented in the 
STM32F405 reference manual, including edge cases for 1, 2, and 3-byte reception where 
ACK/NACK timing differs.

\lstinputlisting[
  language=C,
  caption={\texttt{fsm states encoded}},
  label=code:i2c_state
]{Resources/Code/fsm_states.c}